%% ============================================================================
\probchapter{Random walks and diffusion}{Notes \S1--\S2}

%% ----------------------------------------------------------------------------
\problem{A biased walk}{New}{\S1}

A walker on a line takes independent steps $s_k$, each $+a$ with probability
$p$ and $-a$ with probability $q=1-p$. After $n$ steps it is at
$x_n=\sum_{k=1}^n s_k$.
\begin{parts}
  \item Find the mean and the variance of a single step.
  \item Find $\avg{x_n}$ and the variance of $x_n$.
  \item The drift $\avg{x_n}$ grows like $n$, the spread like $\sqrt n$.
        For $p=0.6$, after how many steps does the drift exceed the
        standard deviation?
\end{parts}

\solution

\solpart{a}{One step}
\begin{equation}
  \avg{s}=a\,p-a\,q=a(2p-1),\qquad \avg{s^2}=a^2 ,
\end{equation}
so
\begin{equation}
  \operatorname{Var}s=\avg{s^2}-\avg{s}^2=a^2\bigl[1-(2p-1)^2\bigr]
  =a^2\bigl[(1-2p+1)(1+2p-1)\bigr]=\ans{4a^2pq}.
\end{equation}
Check: $p=\tfrac12$ gives $a^2$, the unbiased walk; $p=1$ gives $0$, a
deterministic walk.

\solpart{b}{After $n$ steps}
Means add for any steps; variances add because the steps are independent:
\begin{equation}
  \ans{\avg{x_n}=na(2p-1),\qquad \operatorname{Var}x_n=4na^2pq}.
\end{equation}

\solpart{c}{Drift against spread}
The drift exceeds the standard deviation when
\begin{equation}
  na\,\lvert 2p-1\rvert>2a\sqrt{npq}
  \iff n>\frac{4pq}{(2p-1)^2}.
\end{equation}
For $p=0.6$: $4pq=0.96$ and $(2p-1)^2=0.04$, so $n>24$:
\ans{\text{from the 25th step on}}.

%% ----------------------------------------------------------------------------
\problem{How long does diffusion take?}{New}{\S2}

A small molecule in water has $D=1.0\times10^{-9}\,\mathrm{m^2\,s^{-1}}$.
Use the one-dimensional result $\avg{x^2}=2Dt$.
\begin{parts}
  \item How long until the root-mean-square displacement is $1\,\mathrm{mm}$?
  \item How long does it take to cross a cell, $10\,\mu\mathrm{m}$?
  \item Explain the ratio of the two times in one line.
\end{parts}

\solution

\solpart{a}{One millimetre}
\begin{equation}
  t=\frac{\avg{x^2}}{2D}=\frac{(10^{-3}\,\mathrm m)^2}{2\times10^{-9}\,\mathrm{m^2\,s^{-1}}}
  =\ans{500\,\mathrm s}\approx 8\,\mathrm{min}.
\end{equation}

\solpart{b}{Ten micrometres}
\begin{equation}
  t=\frac{(10^{-5}\,\mathrm m)^2}{2\times10^{-9}\,\mathrm{m^2\,s^{-1}}}
  =\ans{0.05\,\mathrm s}.
\end{equation}

\solpart{c}{The ratio}
The time grows as the square of the distance: $100$ times the distance takes
$100^2=10^4$ times as long, and $500\,\mathrm s/0.05\,\mathrm s=10^4$.
